\documentclass[11pt]{article}
\usepackage[utf8]{inputenc}
\usepackage[spanish, es-tabla]{babel}
\usepackage{amsmath, amssymb}
\usepackage{booktabs}
\usepackage{enumitem}
\usepackage[margin=2.5cm]{geometry}

\setlength{\parindent}{0pt}
\setlength{\parskip}{0.65em}

\title{Econometría ICS-294\\Guía de ejercicios: Clase 7}
\author{Francisco Alfaro Medina}
\date{}

\begin{document}
\maketitle

\section*{Objetivo}

Practicar los supuestos clásicos del modelo de regresión lineal múltiple y las propiedades estadísticas del estimador MCO: insesgamiento, varianza y el rol de la exogeneidad.

\section*{Ejercicios}

\begin{enumerate}[label=\textbf{Ejercicio \arabic*.}, leftmargin=*]

\item \textbf{Supuestos clásicos.}

Para cada afirmación, indique qué supuesto clásico está relacionado y explique brevemente.

\begin{enumerate}[label=\alph*)]
    \item La muestra fue seleccionada aleatoriamente desde la población de interés.
    \item Ninguna variable explicativa es combinación lineal exacta de otras.
    \item El valor esperado del error, condicional en \(X\), es cero.
    \item La varianza del error es constante para todos los valores de \(X\).
\end{enumerate}

\item \textbf{Exogeneidad.}

Considere:
\[
salario_i=\beta_0+\beta_1educ_i+\beta_2exper_i+u_i.
\]

\begin{enumerate}[label=\alph*)]
    \item Escriba el supuesto de exogeneidad para este modelo.
    \item Mencione dos razones por las cuales este supuesto podría fallar.
    \item Explique por qué este supuesto es clave para una interpretación causal.
\end{enumerate}

\item \textbf{Sesgo por variable omitida.}

Suponga que habilidad afecta salario y está correlacionada positivamente con educación, pero no se incluye en la regresión.

\begin{enumerate}[label=\alph*)]
    \item ¿Qué supuesto se viola?
    \item ¿Esperaría que el coeficiente de educación esté sesgado hacia arriba o hacia abajo?
    \item Explique la intuición económica.
\end{enumerate}

\item \textbf{Insesgamiento de MCO.}

Use:
\[
\hat{\beta}=(X^\top X)^{-1}X^\top y,
\qquad
y=X\beta+u.
\]

\begin{enumerate}[label=\alph*)]
    \item Muestre que \(\hat{\beta}=\beta+(X^\top X)^{-1}X^\top u\).
    \item Tome esperanza condicional en \(X\).
    \item Explique qué condición permite concluir que \(E(\hat{\beta}\mid X)=\beta\).
\end{enumerate}

\item \textbf{Varianza condicional.}

Bajo homocedasticidad y no autocorrelación:
\[
Var(u\mid X)=\sigma^2 I.
\]

\begin{enumerate}[label=\alph*)]
    \item Escriba \(Var(\hat{\beta}\mid X)\).
    \item Explique qué significa que la matriz sea condicional en \(X\).
    \item ¿Qué ocurre si \(\sigma^2\) aumenta?
\end{enumerate}

\item \textbf{Heterocedasticidad.}

Para cada caso, indique si parece razonable sospechar heterocedasticidad.

\begin{enumerate}[label=\alph*)]
    \item Salario como función de educación y experiencia.
    \item Gasto de hogares como función de ingreso.
    \item Puntaje en una prueba como función de horas de estudio.
\end{enumerate}

Explique en cuál caso esperaría que la varianza del error cambie más con los regresores.

\item \textbf{No multicolinealidad perfecta.}

Explique si hay o no multicolinealidad perfecta en cada caso.

\begin{enumerate}[label=\alph*)]
    \item Se incluye edad y edad al cuadrado.
    \item Se incluye ingreso anual e ingreso mensual multiplicado por 12.
    \item Se incluyen dummies para todas las regiones y una constante.
    \item Se incluye educación y experiencia.
\end{enumerate}

\item \textbf{Matriz de varianzas-covarianzas.}

Suponga que:
\[
\hat{V}(\hat{\beta})=
\begin{bmatrix}
16 & -2 & 1\\
-2 & 4 & 0.5\\
1 & 0.5 & 9
\end{bmatrix}.
\]

\begin{enumerate}[label=\alph*)]
    \item Identifique la varianza de \(\hat{\beta}_1\).
    \item Calcule el error estándar de \(\hat{\beta}_1\).
    \item Identifique la covarianza entre \(\hat{\beta}_1\) y \(\hat{\beta}_2\).
\end{enumerate}

\item \textbf{Discusión conceptual.}

Explique la diferencia entre:
\begin{enumerate}[label=\alph*)]
    \item insesgamiento y precisión;
    \item exogeneidad y homocedasticidad;
    \item sesgo y varianza.
\end{enumerate}

\end{enumerate}

\section*{Preguntas de cierre}

\begin{enumerate}[label=\arabic*.]
    \item ¿Cuál es el supuesto más importante para interpretar causalmente un coeficiente?
    \item ¿Bajo qué supuestos MCO es insesgado?
    \item ¿Qué rol cumple la homocedasticidad en la fórmula de la varianza?
    \item ¿Qué significa que MCO sea BLUE o MELI?
\end{enumerate}

\end{document}
